\section{Discussion}
\label{sec:discussion}

\subsection{Completeness of the Five Primitives}

Section~\ref{sec:completeness-formal} states what we can prove and what we
cannot. Theorem~\ref{thm:partial-complete} establishes partial completeness
under the restriction that the navigation function be computable by a
deterministic Turing machine accessing oracles drawn from five standard
disciplines. Conjecture~\ref{conj:full-complete} states the open question:
that forensic admissibility forces this oracle restriction. We frame the
limits of what is shown: a sixth primitive, were one to be discovered,
would be one whose address-space algebra cannot be simulated under the
oracle disciplines (i)-(v). The candidate cases discussed in
Section~\ref{sec:completeness-formal} (PCAP, blockchain, database tables,
cloud audit logs, mobile extractions, stream-processing state) all reduce.
Quantum-entangled storage is set aside because its non-cloning property
falls outside the deterministic-Turing-machine model that
Theorem~\ref{thm:partial-complete} assumes; whether quantum substrates
admit forensic admissibility in the long run is a separate question.

We are explicit that ``cannot construct a counterexample'' is a weaker
claim than ``no counterexample exists.'' The progression we offer --
oracle-restricted partial completeness theorem, a stated full-completeness
conjecture, and a list of attempted reductions -- is the formal posture
appropriate to the current state of the field. We invite future work to
either prove Conjecture~\ref{conj:full-complete} under a precise
formalization of forensic admissibility, or refute it by exhibiting a
sixth primitive.

\subsection{Hash-Agnosticism and Algorithmic Migration}

Proposition~\ref{prop:c-durability} is parameterized in $\epsilon_H$, the
collision-resistance bound of the hash function. The taxonomy is therefore
hash-agnostic: any hash function with sufficient $\epsilon_H$ admits a
\src{C}-source instantiation, and a deployment may use distinct hash
functions for distinct stores without altering any other component of the
framework. This matters operationally because the cryptographic ground is
shifting. Git is migrating to
SHA-256~\cite{git2020sha256} after demonstration of practical SHA-1
collisions~\cite{stevens2017sha1,leurent2020sha1chosen}. OCI image
descriptors specify SHA-256 (and admit SHA-512 and BLAKE3 as
alternatives)~\cite{ociimagespec_digest}. The Sigstore ecosystem uses
SHA-256 throughout. IPFS uses multihash, allowing future migration. A
forensic implementation must be designed to (a) reject \src{C} stores
whose hash is collision-broken at the deployment date (treating
$\epsilon_H$ as effectively non-negligible), (b) record which hash
function each \src{C}-source result used, so that retrospective
re-evaluation under improved cryptanalytic assumptions remains possible,
and (c) accommodate post-quantum hash candidates as they mature. The
taxonomy itself requires no modification across these transitions, which
is itself an architectural advantage.

\subsection{Threats to Validity}

\paragraph{Construct validity.} Our argument for parser independence treats
the parser as a function $\Bstar \to R$. In practice, parsers may have
side effects (logging, telemetry) and may consume external resources
(allocators, file handles). We have implicitly assumed pure parsers, and
the implementation discipline in Issen reflects this assumption. Parsers
that depend on a specific environment (for example, a parser that requires
network access to look up an external schema) violate the assumption and
should not claim source-agnosticism. The Issen test discipline exercises
parsers in offline mode to validate freedom from network and clock
dependencies, but this discipline is not enforced by the type system; we
view this as a candidate for future work in a Rust effect-system
extension.

\paragraph{Internal validity of Case 3 (parser-independence circularity).}
The empirical demonstration of Theorem~\ref{thm:parser-indep}
(Section~\ref{sec:case3}) uses Issen as the reference implementation, and
Issen was designed to satisfy the theorem. This is a real threat to
internal validity: the implementation was built to satisfy the theory, so
the theory cannot be falsified by it without an independent implementation.
We acknowledge this threat explicitly. A stronger evaluation -- which we
propose as future work -- would re-implement Theorem~\ref{thm:parser-indep}'s
required conditions in a forensic toolkit that was not designed with the
taxonomy in mind (TSK, Volatility, Plaso) and report on the modifications
required to satisfy parser independence. We hypothesize that the
modifications would be substantial but tractable; the magnitude of
required modification is itself a measure of how far from the theorem's
discipline existing tools have drifted.

\paragraph{External validity of Case 1.} The rootkit-concealed cryptominer
scenario is a public-pattern adversarial scenario; the specific incident
was deliberately staged for evaluation. Real incidents may exhibit
adversary capabilities -- kernel-level rootkits, memory-resident loaders
that defeat acquisition tools, supply-chain backdoors -- that the staged
scenario does not. The qualitative conclusion (that fusion across \src{P},
\src{M}, \src{Q} is more defensible than any individual source) is
consistent with practitioner reports~\cite{cohen2011distributed,case2017memory}
but should be evaluated on real incident corpora when available.

\paragraph{External validity of Case 2.} Case 2's hash-pivot demonstration
exercises components of \src{C} that are partially prospective. The OCI and
Sigstore pivots are concrete; the git-reflog inspection is partially
prospective pending completion of \texttt{git-forensic}'s reflog
implementation. We are explicit about this in Section~\ref{sec:case2} and
caution against reading the case as a fully-realized empirical evaluation.

\paragraph{Empirical evaluation does not yet use standard corpora.} A
significant external-validity limitation: we have not yet evaluated Issen
or the parser library against the NIST CFTT
images~\cite{nistcftt} or the
Garfinkel et al.\ standardized forensic corpora~\cite{garfinkel2009corpora}.
These corpora are the appropriate ground truth for a long-lived claim
about parser-source compatibility, and their absence from our evaluation
is a real limitation. We commit to this evaluation as the first item of
future work; we expect the matrix of (parser, source) compatibilities on
the standard corpora to be reported in a companion paper.

\subsection{Threats to Validity}

\paragraph{Construct validity.} Our argument for parser independence treats
the parser as a function $\Bstar \to R$. In practice, parsers may have
side effects (logging, telemetry) and may consume external resources
(allocators, file handles). We have implicitly assumed pure parsers, and
the implementation discipline in Issen reflects this assumption. Parsers
that depend on a specific environment (for example, a parser that requires
network access to look up an external schema) violate the assumption and
should not claim source-agnosticism.

\paragraph{Internal validity of Case 3.} The empirical demonstration of
Theorem~\ref{thm:parser-indep} (Section~\ref{sec:case3}) is a single-host,
single-artifact study. We selected Chrome \texttt{History} because of its
prevalence and its non-trivial schema, but a stronger evaluation would
exercise the parser-independence claim across a battery of artifacts (EVTX,
SRUM, browser histories from multiple vendors, lnk files, prefetch). We
plan such a battery in future work.

\paragraph{External validity of Case 1.} The rootkit-concealed cryptominer
scenario is a public-pattern adversarial scenario; the specific incident
was deliberately staged for evaluation. Real incidents may exhibit
adversary capabilities -- kernel-level rootkits, memory-resident loaders
that defeat acquisition tools, supply-chain backdoors -- that the staged
scenario does not. The qualitative conclusion (that fusion across \src{P},
\src{M}, \src{Q} is more defensible than any individual source) is
consistent with practitioner reports~\cite{cohen2011distributed,case2017memory}
but should be evaluated on real incident corpora when available.

\paragraph{External validity of Case 2.} Case 2's hash-pivot demonstration
exercises components of \src{C} that are partially prospective. The OCI and
Sigstore pivots are concrete; the git-reflog inspection is partially
prospective pending completion of \texttt{git-forensic}'s reflog
implementation. We are explicit about this in Section~\ref{sec:case2} and
caution against reading the case as a fully-realized empirical evaluation.

\subsection{Limitations of the Current Implementation}

Issen as deployed implements \src{P}, \src{M}, \src{L} fully and \src{Q}
for VQL. The \src{C} crate family (\texttt{cas-forensic},
\texttt{git-forensic}, \texttt{sigstore-forensic}) is in active development;
the Rekor and OCI portions of the case study are concrete, but a unified
\src{C} interface that abstracts over git/OCI/Sigstore/IPFS uniformly is
forthcoming. The Pivot engine's trust weighting is currently a single-axis
assignment (higher for \src{C}, lower for \src{P} on disk-resident
artifacts) rather than a fully-developed cost model. We expect to extend
the Pivot engine with a Bayesian trust update analogous to the
hypothesis-update calculus advocated by
Carrier~\cite{carrier2006hypothesis}.

\subsection{Future Work}

Several directions extend this work.

\paragraph{Formal completeness.} A formal completeness theorem -- if one
exists -- would settle the question of whether new evidence sources can
introduce primitives outside our taxonomy. The construction would likely
proceed by characterizing the address-space algebras admissible to
forensic evidence (deterministic enumerable structures with retrieval
semantics) and showing that the five primitives partition the algebra
class.

\paragraph{Empirical evaluation on standard datasets.} The NIST Computer
Forensic Reference Data Sets and DFRWS challenge corpora provide a basis
for an empirical study of parser-independence at scale. We propose to
evaluate Issen's parser library across the NIST CFTT-published images and
report the matrix of (parser, source) compatibility.

\paragraph{Full \src{C} implementation.} \texttt{cas-forensic},
\texttt{git-forensic}, and \texttt{sigstore-forensic} are first priorities.
The interface design must accommodate hash-function diversity (SHA-1
legacy in git, SHA-256 in OCI and Sigstore, multihash in IPFS), and the
trust model must downgrade SHA-1 stores explicitly given known collision
attacks~\cite{stevens2017sha1}.

\paragraph{Source-aware correlation algorithms.} The Pivot engine is a
source-naive cross-product join over normalized timeline events with rule
filtering. Source-aware algorithms -- \eg{}, those that explicitly account
for the differing temporal granularities of \src{M} (instantaneous) versus
\src{L} (linearly ordered) versus \src{Q} (query-batch ordered) versus
\src{C} (no temporal axis at all) -- promise both better recall and better
precision.

\paragraph{Adversarial models.} The attacker-durability propositions
(Propositions~\ref{prop:q-durability} and~\ref{prop:c-durability}) assume
specific adversary capabilities (host-resident, channel-honest). A more
complete model would parameterize the adversary's capability set and ask
which durability claims hold under each. Compromised channels, compromised
investigator storage, and compromised transparency logs all merit explicit
adversarial models.
